\documentclass[10pt,a4paper]{amsart}
\usepackage[margin=19mm]{geometry}
\usepackage{amsmath,amssymb,mathtools}
\usepackage[scr=rsfs,cal=euler]{mathalfa}
\usepackage{xfrac}
\usepackage[unicode,hidelinks,bookmarks=false]{hyperref}
\newcommand{\ie}{i.e.\ }
\newcommand{\cf}{cf.\ }
\newcommand{\resp}{respectively\ }
\newcommand{\Subsection}{\subsection}
\providecommand{\nobreakdash}{\nobreak\hbox{-}}
\setlength{\emergencystretch}{2em}
\multlinegap0pt
\usepackage{bm}
\usepackage{bbm}
\usepackage{dsfont}
\usepackage{xspace}
\usepackage{cancel}
\usepackage{mathtools}
\usepackage{amsthm}
\makeatletter
\@ifundefined{theorem}{\newtheorem{theorem}{Theorem}}{}
\@ifundefined{proposition}{\newtheorem{proposition}[theorem]{Proposition}}{}
\makeatother
\newcommand\magicwand{\mathrel{-\mkern-6mu*}}
\newcommand{\BBI}{\textbf{BBI}\xspace}
\newcommand{\SFour}{\textbf{S4}\xspace}
\title[Embedding Modal Logics into Logics of Bunched Implications]{Embedding Modal Logics into Logics of Bunched Implications}
\author{Daniele Sansoni, Ranald Clouston}
\date{Source: arXiv:2608.07203v1 (CC BY 4.0)}
\begin{document}
\maketitle

\noindent\emph{Source and licence.} Embedding Modal Logics into Logics of Bunched Implications. Version arXiv:2608.07203v1 (CC BY 4.0), distributed under the Creative Commons Attribution 4.0 International licence, as recorded in the arXiv licence metadata for this exact version and shown on the abstract page https://arxiv.org/abs/2608.07203v1. Full rights notice: licenses/arxiv-2608.07203-CC-BY-4.0.txt.\par\smallskip

\noindent\textit{Editorial scope.} Counted unit: the Proposition whose source label is \texttt{Prop:rev\_sound} in the article (source0.tex lines 330--333, source label \texttt{Prop:rev\_sound}), delivered together with the article's own complete original proof of it (source0.tex lines 366--413, one \texttt{proof} environment of 48 lines). No mathematics is written, repaired or re-derived: statement and proof are the article's own text line for line. The proof is detached from the statement in the source: the article states the result at lines 330--333 and prints its proof at lines 366--413, and the proof environment's own header names the statement, so the association is the article's own. Required context retained uncounted: none. Every definition, display and label of those lines is retained unchanged. Omitted from this package and not redistributed: 1--329, 334--365, 414--1469. Nothing inside a retained excerpt is omitted or altered: every retained body is delivered exactly as the article's lines are written, so no omission receipt of any kind accompanies this package. Citations: the retained excerpts carry no citation command at all, so no bibliography entry of the article is transcribed and no reference is redistributed. Reference adaptation: none was needed and none was made; every reference inside the delivered unit resolves inside it, so no unresolved reference or citation is produced. Printed slips kept literal: no printed word, symbol, hypothesis or number of the counted statement or of its proof was altered. Layout and class: the article's own class is \texttt{article} with packages including babel, hyperref, thm-restate, authblk, amsthm, thmtools, amsmath, amssymb, bussproofs, cancel, url, latexsym, multirow, graphicx; the wrapper is the lane's pinned \texttt{amsart} editorial wrapper at 10pt on A4, which restates the theorem-like environments the retained text needs and adds the article's own macro definitions that the retained text needs (\texttt{\textbackslash BBI}, \texttt{\textbackslash SFour}, \texttt{\textbackslash magicwand}), with extra wrapper packages (bm, bbm, dsfont, xspace, cancel, mathtools). The delivered numbering is the unit's own. Assets: no retained span contains a figure environment, a graphics-inclusion command, a PDF-page inclusion, a table or a TikZ picture; no image file is copied into this package, the package declares no asset dependency and no image file is redistributed with it. Licence: version arXiv:2608.07203v1 of this article is distributed under the Creative Commons Attribution 4.0 International licence, as recorded in the arXiv licence metadata for this exact version and shown on the abstract page https://arxiv.org/abs/2608.07203v1. A full rights notice is delivered at licenses/arxiv-2608.07203-CC-BY-4.0.txt. Characters: every retained excerpt is pure ASCII, so the delivered engine, XeLaTeX with its default Latin Modern text font, reports no missing character.
\begin{proposition}[Reverse Translation is Sound]
\label{Prop:rev_sound}
For every formula $\alpha$ of \BBI, if $\vdash_{\BBI} \alpha$ then $\vdash_{\SFour} [\alpha]'$.
\end{proposition}

\begin{proof}[Proof of Proposition~\ref{Prop:rev_sound}]
    By induction on the derivation in \BBI. First, axioms of \BBI translate to formulae provable in \SFour. Indeed, as axioms do not contain $\magicwand$ these axioms translate into the language of propositional logic, and are clearly theorems of classical propositional logic:

  \begin{description}

   \item[1.] $[\alpha \rightarrow I*\alpha]' \;=\; [\alpha]' \rightarrow \top\land[\alpha]'$.

    \item[2.] $[I*\alpha \rightarrow \alpha]' \;=\; \top \land [\alpha]' \rightarrow [\alpha]'$.

    \item[3.] $[\alpha*\beta\rightarrow\beta*\alpha]' \;=\; [\alpha]' \land [\beta]'\rightarrow[\beta]' \land [\alpha]'$.

    \item[4.] $[\alpha*(\beta*\gamma) \rightarrow (\alpha *\beta) * \gamma)]' \;=\; [\alpha]' \land ([\beta]' \land [\gamma]') \rightarrow ([\alpha]' \land [\beta]') \land [\gamma]'$.



  \end{description}

    As the induction step, we see that the translations of inference rules of \BBI preserve derivability in \SFour:

\begin{description}

    \item[1.] Rule R1 (\textit{modus ponens}) is shared by \BBI and \SFour.

\item[2.] The translations of rules R2 and R3 correspond to asking that $(\alpha \land \beta)\rightarrow\gamma$ if and only if $\alpha \rightarrow (\beta \rightarrow (\Box \gamma \lor \gamma))$. In one direction:
\begin{description}
    \item[1.] $\vdash_{\SFour}\alpha \land \beta \rightarrow \gamma$.\textit{ Hyp.}
    \item[2.] $\vdash_{\SFour}(\alpha \land \beta \rightarrow \gamma) \rightarrow (\alpha \land \beta \rightarrow \gamma \lor \Box \gamma)$. \textit{Taut.}
    \item[3.] $\vdash_{\SFour}\alpha \land \beta \rightarrow \gamma \lor \Box \gamma$. \textit{[MP] on 1. and 2.}
    \item[4.] $\vdash_{\SFour}(\alpha \land \beta \rightarrow \gamma \lor \Box \gamma) \rightarrow (\alpha \rightarrow (\beta \rightarrow \Box \gamma \lor \gamma))$. \textit{Taut.}
    \item[5.] $\vdash_{\SFour}\alpha \rightarrow (\beta \rightarrow \Box \gamma \lor \gamma)$. \textit{[MP] on 3. and 4.}
    
\end{description}
In the other direction:
\begin{description}
    \item[1.] $\vdash_{\SFour}\alpha \rightarrow(\beta\rightarrow\Box \gamma \lor \gamma)$. \textit{ Hyp.}
    \item[2.] $\vdash_{\SFour}\Box \gamma \rightarrow \gamma$. \textit{ Modal Axiom Scheme \SFour.}
    \item[3.] $\vdash_{\SFour}(\alpha \rightarrow(\beta\rightarrow\gamma \lor \gamma).$ \textit{Classical Logic on 1. and 2.}
    \item[4.] $\vdash_{\SFour}\alpha \rightarrow(\beta\rightarrow \gamma).$ \textit{Classical Logic on 3.}
    \item[5.] $\vdash_{\SFour}((\alpha \rightarrow (\beta \rightarrow \gamma)) \rightarrow (\alpha \land \beta \rightarrow \gamma)$. \textit{Taut.}
    \item[6.] $\vdash_{\SFour}(\alpha \land \beta \rightarrow \gamma).$ \textit{[MP] on 4. and 5.}
\end{description}

    \item[3.] Translation of rule R4 states that from $(\alpha \rightarrow \gamma)$ and $(\beta \rightarrow \delta)$, we can conclude $( \alpha \land \beta)\rightarrow (\gamma \land \delta)$, which follows via classical propositional logic.

\end{description}
    

\end{proof}

\end{document}
